% graphicx loaded via preamble (includegraphics used below)
% ==============================================================================
% T0 Theory: Document 164 (English)
% Title: T0 Theory and 't Hooft's Cellular Automaton Interpretation of
%        Quantum Mechanics — Determinism from Two Directions
% Author: Johann Pascher
% Date: February 2026
% Reference: G. 't Hooft, The Cellular Automaton Interpretation of Quantum
%            Mechanics, Springer 2016, DOI 10.1007/978-3-319-41285-6
% ==============================================================================

\section*{Abstract}

Gerard 't Hooft's \emph{The Cellular Automaton Interpretation of Quantum Mechanics}
(Springer, 2016) and T0 theory pursue the same fundamental goal: to understand
quantum mechanics as deterministic and to replace the Copenhagen interpretation
with a deeper classical mechanism. This document systematically analyses where
both approaches agree (determinism, rejection of wavefunction collapse, discrete
time structure at the Planck scale, Hilbert space as a tool) and where T0 is more
fundamental: 't Hooft postulates the cellular automaton as a new foundation without
explaining why precisely this structure. T0 derives the discrete spacetime structure
from a single geometric parameter $\xipar = 4/30\,000$. The central thesis of this
document: 't Hooft's cellular automaton is not an independent postulate but a
limiting case of T0 geometry at the sub-Planck scale $L_0 = \xipar \cdot \ell_P$.
It is further shown that 't Hooft's open problems (locality, fermions, gravity,
free parameters of the Standard Model) receive natural answers in T0 through
the time-mass duality $T \cdot m = 1$ and fractal spacetime geometry.

\clearpage

% ==============================================================================

\begin{figure}[htbp]
  \centering
  \includegraphics[width=0.92\textwidth]{../fig/fig1_hierarchy_en}
  \caption{Three layers of quantum interpretation: Copenhagen (probabilistic), CAI by 't~Hooft (automaton postulated), T0 theory (determinism derived from $\xi$). Each level subsumes the previous one.}
  \label{fig:hierarchy_en}
\end{figure}

\section{Introduction: Two Nobel-Prize Intuitions, One Goal}
\label{sec:intro}
% ==============================================================================

\subsection{The Shared Discomfort}

Gerard 't Hooft, Nobel Laureate 1999 for elucidating the quantum structure of
electroweak interactions, has spent recent decades intensely occupied with a
question most of his colleagues considered settled: is quantum mechanics
fundamentally probabilistic, or does a deterministic structure underlie it?

His answer, laid out in detail in \emph{The Cellular Automaton Interpretation of
Quantum Mechanics} (2016), is unambiguous: quantum mechanics is a tool, not a
theory. Beneath it lies a classical-deterministic mechanism operating at the
Planck scale whose effective description yields familiar quantum mechanics.

T0 theory (time-mass duality, FFGFT) arrives at the same conclusion from a
completely different direction. Starting from the observation that all physical
constants emerge from a single geometric parameter $\xipar = 4/30\,000$,
not only does determinism follow as a consequence, but also a concrete geometric
structure of the Planck scale that contains 't Hooft's cellular automaton as a
special case.

\subsection{Structure of This Document}

\begin{enumerate}
    \item \textbf{Section 2}: Core points of 't Hooft's CAI — what it claims,
    what it proves, what it leaves open.
    \item \textbf{Section 3}: Six points of deep agreement between CAI and T0.
    \item \textbf{Section 4}: Four fundamental differences — where T0 goes deeper.
    \item \textbf{Section 5}: The central thesis: the cellular automaton as a
    limiting case of T0.
    \item \textbf{Section 6}: 't Hooft's open problems and T0 answers.
    \item \textbf{Section 7}: Testable predictions from the comparison.
\end{enumerate}

\clearpage

% ==============================================================================
\section{Core Points of the Cellular Automaton Interpretation}
\label{sec:cai_summary}
% ==============================================================================

\subsection{The Fundamental Thesis}

't Hooft begins with a fundamental critique of the Copenhagen interpretation.
The standard interpretation of quantum mechanics leaves several uncomfortable
gaps: wavepacket collapse upon measurement, the non-locality of entanglement,
the impossibility of a classical description of reality, and the apparent
incompatibility with relativity theory at a fundamental level.

't Hooft's counter-proposal in one sentence: \emph{quantum mechanics is the
effective theory of a classical-deterministic system operating at the Planck
scale as a cellular automaton.}

\subsection{The Cogwheel Model as Prototype}

The simplest model of the CAI is the \emph{cogwheel model}: a system with $N$
discrete states $|n\rangle$, $n = 0, \ldots, N-1$, evolving according to a
fixed rule: at each time step $\delta t$, one has $n \to n+1 \pmod{N}$.
This is completely classical-deterministic.

't Hooft now shows that a Hilbert space and Hamiltonian operator can be assigned
to this system:
\begin{equation}
    H_{\text{op}} |n\rangle = \frac{2\pi k}{N \delta t} |n\rangle
    \label{eq:cogwheel_hamilton}
\end{equation}
where $k$ is the energy quantum number. The states obey the Schrödinger equation
— not because the system is quantum mechanical, but because quantum mechanics
\emph{is precisely this}: an effective description of periodic deterministic systems.


\begin{figure}[htbp]
  \centering
  \includegraphics[width=\textwidth]{../fig/fig2_cogwheel_torus_en}
  \caption{Left: CAI cogwheel with $N=8$ discrete states (postulated). Right: T0 time field on geometric torus, lattice spacing $L_0 = \xi \cdot \ell_P$ (derived).}
  \label{fig:cogwheel_en}
\end{figure}

\subsection{Ontological vs.\ Template States}

A central concept of the CAI is the distinction between:
\begin{itemize}
    \item \textbf{Ontological states} $|n\rangle_{\text{ont}}$: the actual states
    of the classical-deterministic system. At each moment, the system is in exactly
    one such state.
    \item \textbf{Template states} $|\psi\rangle$: superpositions of ontological
    states. These are not a picture of reality but description tools for an observer
    who does not know the exact initial conditions.
\end{itemize}

Probabilism in quantum mechanics is, according to 't Hooft, \emph{epistemic}: it
reflects the observer's ignorance of the precise ontological initial conditions,
not a fundamental indeterminacy of nature.

\subsection{Bell's Theorem and Superdeterminism}

The apparently strongest argument against hidden-variable theories is Bell's
theorem: no local, realistic theory can reproduce the quantum correlations.
't Hooft responds with the concept of \emph{superdeterminism}: if the universe
is completely deterministic, then even the settings of measurement apparatuses
are determined by the initial conditions. The degrees of freedom assumed to be
independent in Bell's argument are not.

't Hooft shows that this objection is in principle surmountable, but concedes
that a complete quantitative theory of superdeterminism is still outstanding.

\subsection{The Open Problems}

't Hooft is explicitly transparent about the limits of his approach. At the end
of the book he names explicitly:
\begin{itemize}
    \item The \textbf{locality problem}: how are quantum correlations reproduced
    locally from the automaton?
    \item The \textbf{fermion problem}: how do half-integer spins arise from a
    boson-like automaton?
    \item The \textbf{gravity problem}: the automaton has not been connected to
    general relativity.
    \item The \textbf{free-parameter problem}: the Standard Model has 20+ free
    parameters. The automaton does not explain why they have these values.
    \item The \textbf{information-loss problem}: black holes and unitarity remain
    unresolved.
\end{itemize}

\begin{keyresult}['t Hooft's Programme]
    The CAI is a programme, not a finished theory. It shows that deterministic
    quantum mechanics is \emph{possible}, but it does not yet provide a complete
    derivation of known physics from a deterministic foundation.
\end{keyresult}

\clearpage

% ==============================================================================
\section{Six Points of Deep Agreement}
\label{sec:agreement}
% ==============================================================================

\subsection{1. Determinism as a Fundamental Principle}

Both approaches reject the fundamental probabilism of the Copenhagen interpretation.
In the CAI, randomness is epistemic — it reflects ignorance of initial conditions.
In T0, randomness is equally epistemic: the time field
$\Tfield = \hbar / (mc^2)$ evolves completely deterministically according to the
T0 energy functional:
\begin{equation}
    E_{T0} = \int |\nabla T(x)|^2 \, d^4x
    \label{eq:t0_energy}
\end{equation}
Apparent quantum uncertainty in T0 is the projection of a deterministic
four-dimensional time field onto the three-dimensional spatial hypersurface of
an observer.

\subsection{2. Rejection of Wavepacket Collapse}

't Hooft: collapse is not a physical process but an information update of the
observer about the ontological state of the system. Before the measurement, the
system was in a definite state — the observer simply did not know it.

T0: identical conclusion from a different foundation. The time field $\Tfield$
does not collapse — it evolves continuously. A measurement is an interaction
of two time fields whose correlation generates the apparent collapse structure.

\subsection{3. Hilbert Space as a Tool, Not Ontology}

't Hooft writes explicitly: \emph{quantum mechanics is viewed as a tool, not as
a theory.} Hilbert space does not describe a fundamental reality — it is an
effective description tool for systems about whose microscopic details one has
no complete information.

T0 fully agrees: Hilbert-space states $|\psi\rangle$ are in T0 a compressed
representation of the time-field configuration $T(x,t)$. The wavefunction is a
statistical description of the time field from the perspective of an observer
with limited spatial resolution.

\subsection{4. Discrete Time Structure at the Planck Scale}

't Hooft: the automaton operates with a fundamental time step $\delta t$ of the
order of the Planck time $t_P \approx 5.39 \times 10^{-44}$ s. Below this scale,
no dynamics are defined.

T0: the sub-Planck time scale is not $t_P$ but
\begin{equation}
    T_0 = \xipar \cdot t_P = \frac{4}{30\,000} \cdot t_P
    \approx 7.19 \times 10^{-48} \, \text{s}
    \label{eq:t0_time}
\end{equation}
Both approaches share the conviction that time is discrete at the smallest
scale. T0 provides the exact value of this discreteness — 't Hooft leaves it open.

\subsection{5. Quantum Mechanics as a Limiting Case of Classical Dynamics}

't Hooft shows with concrete models (cogwheel, harmonic rotator, free fields)
that classical-deterministic systems obey Schrödinger equations when one chooses
the appropriate basis. Quantum mechanics \emph{emerges} from classical mechanics.

T0 makes the same statement more precisely: the Schrödinger equation is the
linear approximation of the nonlinear T0 field equation
\begin{equation}
    \nabla^2 T - \frac{1}{c^2}\frac{\partial^2 T}{\partial t^2}
    = -\frac{\xipar}{T^2} \cdot |\nabla T|^2
    \label{eq:t0_field_equation}
\end{equation}
for weak time-field gradients, i.e.\ for $|\nabla T| \ll 1/\xipar$.

\subsection{6. Arrow of Time as an Emergent Phenomenon}

't Hooft: the arrow of time can be explained from the deterministic automaton
if one introduces information loss at collisions — certain microstates collapse
onto the same macrostate.

T0: the arrow of time follows geometrically from $T \cdot m = 1$. Since $T > 0$
is algebraically enforced and torus winding defines a topologically protected
direction, the arrow of time is not a thermodynamic approximation but a
fundamental geometric fact (cf.\ Document 157).

\begin{discoveryBox}[Deep Convergence]
    Two independently developed theories — 't Hooft's CAI from the perspective
    of quantum field theory and T0 from the perspective of geometric unification —
    converge on six fundamental points. This is not coincidence: both push against
    the same wall of the standard interpretation and find similar breakthroughs.
\end{discoveryBox}

\clearpage

% ==============================================================================
\section{Four Fundamental Differences: Where T0 Goes Deeper}
\label{sec:differences}
% ==============================================================================

\subsection{Difference 1: Postulate vs.\ Derivation}

\paragraph{'t Hooft:} The cellular automaton is \emph{postulated} as a new
fundamental structure. It is assumed that nature at the Planck scale functions
like an automaton with discrete states and a deterministic evolution rule.
Why \emph{this} structure? The book provides no answer.

\paragraph{T0:} The discrete spacetime structure is \emph{derived} from
$\xipar = 4/30\,000$. The parameter $\xipar$ is not free — it emerges from the
geometry of three-dimensional space at the Planck scale. The cell sizes of the
automaton are thereby fixed:
\begin{equation}
    \ell_{\text{cell}} = L_0 = \xipar \cdot \ell_P, \quad
    t_{\text{step}} = T_0 = \xipar \cdot t_P
    \label{eq:cell_size}
\end{equation}

\begin{table}[h]
    \centering
    \caption{Postulate vs.\ derivation: fundamental structures}
    \label{tab:postulate_vs_derivation}
    \resizebox{\textwidth}{!}{\begin{tabular}{@{}lll@{}}
        \toprule
        Structure & 't Hooft CAI & T0 Theory \\
        \midrule
        Discrete time structure & Postulated ($\delta t \sim t_P$) & Derived ($T_0 = \xipar \cdot t_P$) \\
        Cell size & Open ($\sim \ell_P$) & Fixed ($L_0 = \xipar \cdot \ell_P$) \\
        Evolution rule & Postulated (automaton) & Derived from $E_{T0}$ \\
        Free parameters & Several & Zero \\
        Why this value? & No answer & $\xipar$ from geometry \\
        \bottomrule
    \end{tabular}}
\end{table}

\subsection{Difference 2: Free Parameters of the Standard Model}

\paragraph{'t Hooft:} The book explicitly acknowledges that the CAI does not
explain the 20+ free parameters of the Standard Model. The quark masses, coupling
constants, the Weinberg angle — all remain unexplained.

\paragraph{T0:} All these parameters emerge from $\xipar$. Particle masses follow
from $E_n = E_0 / \xipar^n$ with $E_0 = m_e c^2$:
\begin{equation}
    m_n = m_e \cdot \xipar^{-n}, \quad n = 0, 1, 2, \ldots
    \label{eq:mass_spectrum}
\end{equation}
This formula reproduces 98\% of all known particle masses. The fine structure
constant follows from $\alpha = \xipar \cdot E_0^2$.

\subsection{Difference 3: Gravity and Spacetime}

\paragraph{'t Hooft:} The integration of gravity into the CAI is explicitly listed
as an open problem. The book contains an appendix on gravity in 2+1 dimensions
as a test case, but no complete 3+1-dimensional theory.

\paragraph{T0:} Gravity emerges naturally from the time-field gradient:
\begin{equation}
    \vec{g}(x) = -c^2 \nabla \ln T(x) = \frac{c^2}{T} \nabla T
    \label{eq:t0_gravity}
\end{equation}
Newton's gravitational potential $\Phi = -GM/r$ corresponds to
$\ln T \propto GM/(rc^2)$, and general relativity emerges as the classical
limiting case of the T0 field equations. The gravitational constant $G$ also
emerges from $\xipar$:
\begin{equation}
    G = \frac{\xipar^2 \cdot \ell_P^2 \cdot c^3}{\hbar}
    \label{eq:G_from_xi}
\end{equation}

\subsection{Difference 4: The Fermion Problem}

\paragraph{'t Hooft:} Fermions (half-integer spin, Pauli principle) are
particularly difficult to obtain from a boson-like automaton. The book contains
a chapter on fermions but itself explains that the construction is still incomplete.

\paragraph{T0:} Fermions arise naturally from the topology of the time field.
The binding condition $T \cdot m = 1$ combined with the toroidal topology of the
T0 universe generates half-integer winding numbers corresponding to fermionic
spin. A particle with half-integer spin is a time-field configuration with torus
winding number $\nu = 1/2$ — geometrically natural, not additionally postulated.

\clearpage

% ==============================================================================
\section{The Cellular Automaton as a Limiting Case of T0}
\label{sec:ca_as_limit}
% ==============================================================================

\subsection{The Central Thesis}

We now formulate the main claim of this document precisely:

\begin{revolutionary}[Main Thesis: Inclusion]
    't Hooft's cellular automaton is not an independent foundation of physics
    but an effective limiting case of T0 geometry at the sub-Planck scale.
    Concretely: the CAI describes the discrete time evolution of the T0 time
    field $T(x,t)$ in the lattice approximation $\Delta x = L_0$, $\Delta t = T_0$.
\end{revolutionary}


\begin{figure}[htbp]
  \centering
  \includegraphics[width=0.95\textwidth]{../fig/fig3_inclusion_en}
  \caption{Inclusion diagram: the CAI (orange) is contained as a limiting case within T0 theory (green). Outer green elements are T0 contributions beyond the CAI.}
  \label{fig:inclusion_en}
\end{figure}

\subsection{The Connection in Detail}

In the CAI, the time evolution of the universe is described by a unitary evolution
matrix $U_{\text{op}}(\delta t) = e^{-iH_{\text{op}}\delta t}$ acting on
ontological states. The Hamiltonian $H_{\text{op}}$ is formally defined but not
derived from deeper principles.

In T0, the equivalent is the time-field energy functional $E_{T0}$.
The time evolution of the discrete T0 lattice at scale $L_0$ is:
\begin{equation}
    T(x + L_0, t + T_0) = T(x, t) - \xipar \cdot T_0 \cdot
    \frac{\partial E_{T0}}{\partial T(x,t)}
    \label{eq:t0_discrete_evolution}
\end{equation}

This is exactly the structure of a cellular automaton:
\begin{itemize}
    \item Each cell $(x, t)$ contains the value $T(x,t)$ — this is the
    \emph{ontological state} in 't Hooft's sense.
    \item The evolution rule is local and deterministic — exactly as in the CAI.
    \item The CAI Hamiltonian $H_{\text{op}}$ corresponds in T0 to the second
    variation of the energy functional $\delta^2 E_{T0}$.
\end{itemize}

\subsection{Why T0 Delivers More}

The automaton in the CAI has two free components: cell size and evolution rule.
T0 fixes both:
\begin{align}
    \ell_{\text{cell}} &= L_0 = \xipar \cdot \ell_P
    \quad \Longleftarrow \quad \xipar = \frac{4}{30\,000}
    \label{eq:cell_size_fixed} \\
    \text{Evolution rule} &: T \to T - \xipar T_0 \frac{\delta E_{T0}}{\delta T}
    \quad \Longleftarrow \quad E_{T0} = \int |\nabla T|^2 d^4x
    \label{eq:rule_fixed}
\end{align}

The automaton is not freely choosable — it is geometrically determined.

\subsection{The Cogwheel Model in T0 Language}

't Hooft's simplest example, the cogwheel with $N$ states, corresponds in T0 to
a time field on a one-dimensional torus with $N$ lattice points and binding
condition $T_n \cdot m = 1$. The periodic evolution rule $n \to n + 1 \pmod{N}$
arises from the geometric periodicity of the torus. The energy eigenvalues from
\eqref{eq:cogwheel_hamilton} correspond to the harmonic frequencies
$f_k = k/(N \cdot T_0)$ of the T0 time field on the torus.

\begin{insightBox}[Geometry Replaces Postulate]
    The cogwheel model is not an \emph{example} of a possible automaton — it is
    the \emph{only} model compatible with T0 torus geometry. The periodic
    structure is not chosen; it is geometrically enforced.
\end{insightBox}

\clearpage

% ==============================================================================
\section{'t Hooft's Open Problems and T0 Answers}
\label{sec:open_problems}
% ==============================================================================


\begin{figure}[htbp]
  \centering
  \includegraphics[width=0.62\textwidth]{../fig/fig4_radar_en}
  \caption{Radar chart of open problems. Orange: CAI ('t~Hooft). Green: T0 theory. Score: 0\,=\,open, 2\,=\,resolved.}
  \label{fig:radar_en}
\end{figure}

\subsection{The Locality Problem}

\paragraph{CAI problem:} Quantum entanglement produces non-local correlations
apparently difficult to reconcile with a locally operating automaton rule.
't Hooft argues for superdeterminism but concedes that a quantitative theory
is missing.

\paragraph{T0 answer:} In T0, entangled particles are connected through a shared
time-field configuration $T_{\text{AB}}(x_A, x_B, t)$. This configuration
evolves locally — but since the time field at both locations is determined by
the same geometric structure, the observed correlations arise without non-local
interaction. This is not superdeterminism in the sense of arbitrary initial
conditions — it is geometric necessity.

\subsection{The Fermion Problem}

\paragraph{CAI problem:} Half-integer spins do not arise naturally from an
integer automaton. The chapter on fermions in the book remains unfinished.

\paragraph{T0 answer:} The toroidal topology of the T0 universe allows time-field
configurations with half-integer winding numbers. A fermion is mathematically a
time-field configuration with
\begin{equation}
    \oint_\gamma \frac{d T}{T} = \pi \quad \text{(half-integer winding)}
    \label{eq:fermion_winding}
\end{equation}
This is geometrically stable (topological protection) and requires no additional
postulate.

\subsection{The Gravity Problem}

\paragraph{CAI problem:} No complete connection to general relativity. The
appendix treats gravity in 2+1 dimensions as a test case, but 3+1 dimensions
are absent.

\paragraph{T0 answer:} Equation \eqref{eq:t0_gravity} directly connects the
time field to gravitational acceleration. In the weak-field approximation, T0
reproduces Einstein's general relativity. In the strong-field approximation
(near singularities), T0 gives corrections of order $\xipar^2$.

\subsection{The Free-Parameter Problem}

\paragraph{CAI problem:} The 20+ free parameters of the Standard Model are not
explained. The book mentions this problem but describes it as outside the current
framework.

\paragraph{T0 answer:} All parameters follow from $\xipar$:
\begin{table}[h]
    \centering
    \caption{Free parameters of the Standard Model in T0}
    \label{tab:sm_parameters}
    \resizebox{\textwidth}{!}{\begin{tabular}{@{}lll@{}}
        \toprule
        Parameter & Standard Model & T0 \\
        \midrule
        Fine structure constant $\alpha$ & Measured: $1/137.036$ & $\alpha = \xipar \cdot E_0^2$ \\
        Gravitational constant $G$ & Measured & $G = \xipar^2 \ell_P^2 c^3/\hbar$ \\
        Electron mass $m_e$ & Measured & $m_e = E_0/c^2$, $E_0 = \hbar c/(\xipar \cdot \ell_P)$ \\
        Muon mass $m_\mu$ & Measured & $m_\mu = m_e/\xipar$ \\
        Tau mass $m_\tau$ & Measured & $m_\tau = m_e/\xipar^2$ \\
        Planck constant $h$ & Measured & Emerges from $\xipar$ \\
        All others & 15+ free & From $\xipar$ \\
        \bottomrule
    \end{tabular}}
\end{table}

\subsection{The Information-Loss Problem}

\paragraph{CAI problem:} 't Hooft argued in earlier work that information may be
lost during black hole formation — in conflict with the unitarity of quantum
mechanics. The book addresses this without providing a complete resolution.

\paragraph{T0 answer:} In the T0 framework, information is never lost. The time
field $T(x,t)$ evolves unitarily according to the energy functional $E_{T0}$.
Black holes in T0 are regions of extreme time-field compression.
Crucially, the fractal correction $K_{\text{frak}}$ imposes a hard geometric
lower bound. The time field \emph{cannot} be compressed to zero — a minimal
$T$ exists:
\begin{equation}
    T \geq T_{\min} = \xi \cdot t_P \approx 7.19 \times 10^{-48}\,\text{s}
    \quad \Longrightarrow \quad m \leq m_{\max} = \frac{\hbar}{T_{\min}\,c^2}
    \label{eq:t_min_bh}
\end{equation}
Black holes in T0 are therefore not true singularities — the time field
compresses to $T_{\min}$ but remains finite and retains well-defined gradients.
The information paradox dissolves: the time field evolves unitarily, carries
full information, and no divergence ever occurs.

\clearpage

% ==============================================================================
\section{Testable Predictions from the Comparison}
\label{sec:predictions}
% ==============================================================================


\begin{figure}[htbp]
  \centering
  \includegraphics[width=\textwidth]{../fig/fig5_scales_en}
  \caption{Scale hierarchy from human scale to the T0 sub-Planck scale. The T0 lattice constant $L_0 = \xi \cdot \ell_P$ lies by factor $\xi \approx 1.33 \times 10^{-4}$ below the Planck length.}
  \label{fig:scales_en}
\end{figure}

\subsection{Prediction 1: Cell-Size Signature}

't Hooft leaves the cell size $\delta t$ open (of order $t_P$). T0 fixes it to
$T_0 = \xipar \cdot t_P$. The ratio
\begin{equation}
    \frac{T_0}{t_P} = \xipar = \frac{4}{30\,000} \approx 1.33 \times 10^{-4}
    \label{eq:cell_ratio}
\end{equation}
is in principle measurable: in precision experiments at the boundary of Planck
resolution, the discrete time structure should become visible at $T_0$, not at
$t_P$.

\subsection{Prediction 2: Automaton States and Particle Masses}

If the cellular automaton is a limiting case of T0, then the energy eigenvalues
of the automaton must exactly correspond to the T0 mass spectrum:
\begin{equation}
    E_n^{\text{automaton}} = E_0 \cdot \xipar^{-n}, \quad n = 0, 1, 2, \ldots
    \label{eq:automaton_mass}
\end{equation}
This is a direct connection between the CAI structure and experimentally
measured particle masses.

\subsection{Prediction 3: Bell Correlations and $\xipar$ Corrections}

In the CAI, superdeterminism is not worked out quantitatively. T0 provides a
concrete prediction: Bell correlations should deviate from ideal quantum
predictions on the scale $\xipar \approx 10^{-4}$:
\begin{equation}
    \langle AB \rangle_{T0} = \langle AB \rangle_{\text{QM}}
    \cdot \left(1 - \xipar \cdot \frac{\Delta x}{\ell_P}\right)
    \label{eq:bell_correction}
\end{equation}
where $\Delta x$ is the spatial separation of the measurement locations.
This correction is just below the current detection threshold but is in
principle testable.

\subsection{Prediction 4: Hamiltonian from T0}

In the CAI, the Hamiltonian remains a tool without deep origin. T0 predicts
that $H_{\text{op}}$ corresponds to the second variation of the energy
functional:
\begin{equation}
    H_{\text{op}} = \frac{\delta^2 E_{T0}}{\delta T^2}
    \bigg|_{T = T_{\text{vac}}}
    \label{eq:hamilton_from_t0}
\end{equation}
This is a testable statement: the spectra of quantum systems should be calculable
from the T0 time-field energy functional — without free parameters.

\clearpage

% ==============================================================================

\begin{figure}[htbp]
  \centering
  \includegraphics[width=\textwidth]{../fig/fig6_agreements_en}
  \caption{Agreements (green, left) and fundamental differences (blue, right) between the CAI and T0 theory.}
  \label{fig:agreements_en}
\end{figure}

\section{Full Comparison: CAI and T0}
\label{sec:full_comparison}
% ==============================================================================

\begin{table}[h]
    \centering
    \caption{Systematic comparison: CAI ('t Hooft 2016) and T0 theory}
    \label{tab:full_comparison}
    \resizebox{\textwidth}{!}{\begin{tabular}{@{}lll@{}}
        \toprule
        Aspect & CAI ('t Hooft) & T0 Theory \\
        \midrule
        \textbf{Core claim} & QM is effective theory of a CA & QM emerges from time-field geometry \\
        \textbf{Determinism} & Yes (ontological states) & Yes ($T(x,t)$ deterministic) \\
        \textbf{Probabilism} & Epistemic & Epistemic \\
        \textbf{Wavepacket collapse} & No real collapse & No real collapse \\
        \textbf{Hilbert space} & Tool & Tool (projection of time field) \\
        \textbf{Time discreteness} & Postulated ($\delta t \sim t_P$) & Derived ($T_0 = \xipar \cdot t_P$) \\
        \textbf{Cell structure} & Postulated & Geometrically determined \\
        \textbf{Bell's theorem} & Superdeterminism (qualitative) & Geometric correlations ($\xipar$) \\
        \textbf{Free parameters} & Not explained & Zero (all from $\xipar$) \\
        \textbf{Fermions} & Incomplete & Topological winding numbers \\
        \textbf{Gravity} & Open & From time-field gradient \\
        \textbf{Arrow of time} & Information loss & Geometrically enforced \\
        \textbf{Black holes} & Information paradox open & Time-field compression, unitary \\
        \textbf{Standard Model} & 20+ free parameters open & All from $\xipar$ \\
        \textbf{Status} & Programme, still incomplete & Complete geometric framework \\
        \bottomrule
    \end{tabular}}
\end{table}

\clearpage

% ==============================================================================
\section{Summary}
\label{sec:summary}
% ==============================================================================

Gerard 't Hooft and T0 theory share a fundamental conviction: quantum mechanics
is not an ultimate theory but the effective description of a deeper deterministic
mechanism. Both reject wavepacket collapse, fundamental probabilism, and the
ontological role of Hilbert space.

On six fundamental points, CAI and T0 agree: determinism, rejection of collapse,
Hilbert space as a tool, discrete time structure at the Planck scale, quantum
mechanics as a limiting case of classical dynamics, and the arrow of time as
an emergent phenomenon.

On four points, T0 is more fundamental: instead of postulating the automaton,
T0 derives the discrete spacetime structure from $\xipar = 4/30\,000$. The free
parameters of the Standard Model emerge from $\xipar$. Gravity follows from
the time-field gradient. Fermions arise from topological winding numbers.

The central thesis: 't Hooft's cellular automaton is a limiting case of T0
geometry — the discrete lattice approximation of the time field $T(x,t)$ at
the sub-Planck scale $L_0 = \xipar \cdot \ell_P$. 't Hooft's programme shows
that deterministic quantum mechanics is \emph{possible}. T0 shows why it is
\emph{necessary} and what it concretely looks like.

\begin{philosophical}[Einstein's Last Word]
    Einstein was right: God does not play dice. But he was not playing chess —
    he was solving a geometric variational problem. The minimum of
    $E_{T0} = \int |\nabla T|^2 d^4x$ subject to the constraint $T \cdot m = 1$
    is the universe.
\end{philosophical}

\begin{thebibliography}{99}

\bibitem{thooft_2016}
G.\ 't Hooft,
\newblock \emph{The Cellular Automaton Interpretation of Quantum Mechanics},
\newblock Fundamental Theories of Physics 185,
\newblock Springer, Cham, 2016.
\newblock DOI: 10.1007/978-3-319-41285-6.

\bibitem{thooft_arxiv}
G.\ 't Hooft,
\newblock The Cellular Automaton Interpretation of Quantum Mechanics,
\newblock arXiv:1405.1548 [quant-ph] (2014, rev.\ 2015).

\bibitem{pascher_xi}
J.\ Pascher,
\newblock The Geometric Origin Parameter $\xi = 4/30\,000$,
\newblock T0 Documentation Series, Document 009 (2025).

\bibitem{pascher_fundamentals}
J.\ Pascher,
\newblock T0 Theory: Time-Mass Duality as a Fundamental Principle,
\newblock T0 Documentation Series, Document 003 (2025).

\bibitem{pascher_arrow}
J.\ Pascher,
\newblock Arrow of Time as a Geometric Consequence of Spacetime Structure,
\newblock T0 Documentation Series, Document 157 (2026).

\bibitem{pascher_qm_opt}
J.\ Pascher,
\newblock T0 Quantum Mechanics: Geometric Formalism and Optimisation Strategies,
\newblock T0 Documentation Series, Document 034 (2025).

\bibitem{pascher_ising}
J.\ Pascher,
\newblock T0 Theory and Coherent Ising Machines,
\newblock T0 Documentation Series, Document 161 (2026).

\bibitem{pascher_uv}
J.\ Pascher,
\newblock Ultraviolet Catastrophe, Planck's Quantization and Einstein's Light Quanta,
\newblock T0 Documentation Series, Document 163 (2026).

\bibitem{bell_1964}
J.\ S.\ Bell,
\newblock On the Einstein Podolsky Rosen paradox,
\newblock \emph{Physics Physique Fizika} \textbf{1}, 195--200 (1964).

\bibitem{bohm_1952}
D.\ Bohm,
\newblock A suggested interpretation of the quantum theory in terms of
hidden variables I \& II,
\newblock \emph{Physical Review} \textbf{85}, 166--193 (1952).

\end{thebibliography}
